%%=============================================================================
%% CAPITULO 3 -- PROCEDIMENTOS METODOLOGICOS
%% Texto de esqueleto (provisorio), a ser reescrito pela autora.
%%=============================================================================

\section{Procedimentos metodológicos}

O método da pesquisa parte da constatação de que o controle manual da
conformidade no Sistema Eletrônico de Informações é pouco escalável e sujeito a
falhas diante do grande volume de processos, o que motiva o estudo de soluções
automatizadas de monitoramento contínuo.

A investigação examina como técnicas de inteligência artificial e de
processamento de linguagem natural podem ser aplicadas aos documentos do SEI
para identificar dados pessoais e sensíveis, seguindo o fluxo que vai da leitura
dos documentos até a emissão de alertas de conformidade.

% TODO: referência -- a obra a citar para o Método DELPHI ainda será escolhida
% pela autora.
A validação dos critérios segue o Método DELPHI. O painel é anônimo e reúne de
6 a 20 especialistas, que respondem a até três rodadas de questionário on-line.
Em cada rodada, os especialistas atribuem a cada critério notas de 1 a 4 para a
relevância e para a clareza.

% TODO: referência -- a obra a citar para o Índice de Validade de Conteúdo
% ainda será escolhida pela autora.
O consenso é medido, item a item, pelo Índice de Validade de Conteúdo (IVC),
complementado pela mediana das notas. A partir da segunda rodada, os
especialistas recebem o resultado do grupo antes de avaliar novamente os itens.
Os itens que não alcançam o consenso são retirados do conjunto de critérios.
